\documentclass[11pt,a4paper]{amsart}
\usepackage{amsmath,amssymb,amsthm}
\usepackage[margin=2.6cm]{geometry}
\usepackage{booktabs}
\usepackage[colorlinks=true,linkcolor=blue,citecolor=blue,urlcolor=blue]{hyperref}

\newtheorem{theorem}{Theorem}[section]
\newtheorem{proposition}[theorem]{Proposition}
\newtheorem{lemma}[theorem]{Lemma}
\newtheorem{corollary}[theorem]{Corollary}
\theoremstyle{definition}
\newtheorem{remark}[theorem]{Remark}
\newtheorem{example}[theorem]{Example}

\newcommand{\Gg}{\mathcal{G}}
\newcommand{\Gs}{\mathcal{G}_{\star}}
\newcommand{\R}{\mathbb{R}}
\newcommand{\eps}{\varepsilon}
\DeclareMathOperator{\artanh}{artanh}

\title[A sharp constant for Gurland's ratio]{A sharp localization constant
for Gurland's ratio of the gamma function}

\author{DeepSeek V4.1 Flash}
\author{GPT 5.6 Sol}
\author{Grok 4.7}
\author{Claude Opus 5}
\date{\today}

\subjclass[2020]{Primary 33B15; Secondary 11M35, 26D07, 41A58}
\keywords{Gurland's ratio, gamma function, Hurwitz zeta function, Jensen gap,
trigamma mean, remainder estimates}

\begin{document}

\begin{abstract}
Gurland's ratio $\Gg(x,y)=\Gamma(x)\Gamma(y)/\Gamma^{2}\big(\tfrac{x+y}{2}\big)$ is the
exponential of the Jensen gap of $\log\Gamma$. The gap has the classical expansion
\[
  \log\Gg(x,y)=\sum_{k\ge1}\frac1k\Big(\frac{x-y}{2}\Big)^{2k}
  \zeta\Big(2k,\frac{x+y}{2}\Big),
\]
in the form recorded by Srivastava, and the series converges for all $x,y>0$: its radius
is the distance from the midpoint to the pole of $\Gamma$ at the origin. The same
specialization for the modified ratio $\Gs$ therefore holds on the whole quadrant, and
the restriction $Q<1$ is not required for identifying $\log\Gs$ with its series. That
identification settles two questions posed by Wi\'sniewska (arXiv:2512.07028); it is the
classical expansion at the point $1+\tfrac{x+y}{2}$. We give two-sided remainder
estimates whose geometric rate $\rho=|x-y|/(2+x+y)$ is optimal in the sense that
$\tfrac1m\log R_m\to 2\log\rho$. The prefactors on the two sides may differ by a large
factor. The rate satisfies $\rho\le Q$, with equality only for $x=y$, while $Q$ is
unbounded. For large arguments,
$\log\Gs=(1+v)\,\Lambda(\rho)\,(1+o(1))$, where
$\Lambda(\rho)=2\big(\log 2-H_b\big(\tfrac{1+\rho}{2}\big)\big)$ and $H_b$ is the binary
entropy formed with the same logarithm. This is a leading-order statement. The expansion
extends to any number of variables, with applications to the beta function and to
certified computation. The principal result localizes the mean-value parameter $t(x,y)$
in the two-sided bound of Yang and Zheng. Writing
\[
  \theta(x,y)=\frac{t(x,y)-\sqrt{xy}}{\tfrac{x+y}{2}-\sqrt{xy}},
\]
we prove $\theta(x,y)>\tfrac23$. The constant $\tfrac23$ is optimal and is not attained.
The proof is self-contained. High-precision checks accompany the statements.
\end{abstract}

\maketitle

\section{Introduction}

\subsection{Background}
For $x,y>0$ Gurland's ratio \cite{Gurland56} is
\begin{equation}\label{eq:G}
  \Gg(x,y)=\frac{\Gamma(x)\,\Gamma(y)}{\Gamma^{2}\big(\frac{x+y}{2}\big)},
\end{equation}
and Gurland's inequality states that $\Gg(x,y)>1$ for $x\ne y$. Since $\log\Gamma$ is
strictly convex on $(0,\infty)$ (Bohr--Mollerup), $\log\Gg(x,y)$ is exactly the
Jensen gap
\begin{equation}\label{eq:jensen}
  \log\Gg(x,y)=\log\Gamma(x)+\log\Gamma(y)-2\log\Gamma\Big(\frac{x+y}{2}\Big)\;\ge\;0,
\end{equation}
so Gurland's ratio is the exponential of the second-order Jensen gap of $\log\Gamma$.
The modified ratio
\begin{equation}\label{eq:Gs}
  \Gs(x,y)=\frac{\Gamma(1+x)\,\Gamma(1+y)}{\Gamma^{2}\big(1+\frac{x+y}{2}\big)}
            =\frac{4xy}{(x+y)^{2}}\,\Gg(x,y)
\end{equation}
is the form studied in \cite{Wisniewska25}. The ratio \eqref{eq:G} and its relatives were
studied by Gautschi \cite{Gautschi74}, Kairies \cite{Kairies83}, Merkle \cite{Merkle05},
Alzer \cite{Alzer01} and Chen--Choi \cite{ChenChoi17}; see \cite{LaforgiaNatalini13} for
a survey.

Writing
\begin{equation}\label{eq:uv}
  u:=\frac{x-y}{2},\qquad v:=\frac{x+y}{2},\qquad\text{so that }x=v+u,\quad y=v-u,
\end{equation}
the Weierstrass product yields the finite expansion
\begin{equation}\label{eq:Wthm}
  \log\Gs(x,y)=\sum_{k=1}^{m-1}\frac{1}{k}\Big(\frac{x-y}{2}\Big)^{2k}
     \zeta\Big(2k,1+\frac{x+y}{2}\Big)+R_{m}(x,y),\qquad m\ge2,
\end{equation}
with the remainder estimate
\begin{equation}\label{eq:Wrem}
  0\le R_{m}(x,y)\le
  \frac1m\Big(\frac{x-y}{2}\Big)^{2m}\sum_{n=1}^{\infty}\frac{1}{(n+x)^{m}(n+y)^{m}},
\end{equation}
where $\zeta(s,a)=\sum_{n\ge0}(n+a)^{-s}$ is the Hurwitz zeta function
\cite[Thm.~1]{Wisniewska25} (the remainder is denoted $R^{m}_{n}$ there). From
\eqref{eq:Wthm}--\eqref{eq:Wrem} it is deduced in \cite[Cor.~1]{Wisniewska25} that the
condition
\begin{equation}\label{eq:Q}
  Q(x,y):=\frac{|x-y|}{2\,(1+\sqrt{xy})}<1
\end{equation}
implies $\log\Gs(x,y)=S_{\infty}(x,y)$, where $S_\infty=\lim_{m\to\infty}S_m$ and $S_m$
denotes the partial sum in \eqref{eq:Wthm}. The same source gives the bilateral inequality
\begin{equation}\label{eq:Wbil}
  \Big(\frac{x-y}{2}\Big)^{2}\zeta\Big(2,1+\frac{x+y}{2}\Big)\le\log\Gs(x,y)
  \le\Big(\frac{x-y}{2}\Big)^{2}\zeta\big(2,1+\sqrt{xy}\,\big).
\end{equation}
This two-sided bound is a case of Corollary~11 and Remark~5 of Yang and Zheng
\cite{YangZheng19}. With base point $1$, $p=x$ and $q=y$, those results state that
\begin{equation}\label{eq:YZ}
  \exp\Big[\frac{(p-q)^{2}}{4}\psi'(1+r_{1})\Big]
  \;\le\;\frac{\Gamma(1+p)\,\Gamma(1+q)}{\Gamma^{2}\big(1+\frac{p+q}{2}\big)}
  \;\le\;\exp\Big[\frac{(p-q)^{2}}{4}\psi'(1+r_{2})\Big],
\end{equation}
with the best constants $r_{1}=\frac{p+q}{2}$ and $r_{2}=\sqrt{pq}$ whenever $p,q>0$. As
$\psi'=\zeta(2,\cdot)$ and $\big(\frac{x-y}{2}\big)^{2}=u^{2}$, this is
\eqref{eq:Wbil}. Since $\zeta(2,\cdot)$ is continuous and strictly decreasing,
\eqref{eq:Wbil} produces a unique parameter
$t=t(x,y)\in\big(\sqrt{xy},\tfrac{x+y}{2}\big)$ with
\begin{equation}\label{eq:Wt}
  \log\Gs(x,y)=\Big(\frac{x-y}{2}\Big)^{2}\zeta(2,1+t).
\end{equation}
Thus the bound \eqref{eq:Wbil} and the existence of $t$ are contained in
\cite{YangZheng19}. The question that remains is quantitative: where, inside
$\big(\sqrt{xy},\frac{x+y}{2}\big)$, the point $t$ lies. The ratio
\begin{equation}\label{eq:theta}
  \theta(x,y)
  =\frac{t(x,y)-\sqrt{xy}}{\frac{x+y}{2}-\sqrt{xy}}
\end{equation}
measures that position. For $x\ne y$ one has $\theta(x,y)\in(0,1)$, the endpoint $0$
being the geometric mean and the endpoint $1$ the arithmetic mean. Theorems
\ref{thm:loc} and \ref{thm:sharp} locate $\theta$.

Three questions are posed in \cite[\S2]{Wisniewska25}.

\smallskip
\noindent\textbf{Open problem 1.} Does $R_{m}(x,y)\to0$ hold for every $x,y>0$?

\noindent\textbf{Open problem 2.} Does $\log\Gs(x,y)=S_{\infty}(x,y)$ hold in general, and
if not, how large can the difference be?

\noindent\textbf{Open problem 3.} How is $t(x,y)$ located in the interval
$\big(\sqrt{xy},\tfrac{x+y}{2}\big)$? What is its distance to the midpoint of that interval?
\normalsize

\subsection{Relation to the literature, and what is new here}
The identity in Theorem \ref{thm:main} is a symmetrised Taylor expansion and is
classical. Since $\psi^{(n)}(z)=(-1)^{n+1}n!\,\zeta(n+1,z)$ for
$\operatorname{Re}z>0$ and $n\ge1$ \cite[\S5.15]{DLMF}, the Taylor series of $\log\Gamma$
about $a>0$ is
\begin{equation}\label{eq:taylorlog}
  \log\Gamma(a+w)=\log\Gamma(a)+\psi(a)\,w
  +\sum_{n\ge2}\frac{(-1)^{n}}{n}\zeta(n,a)\,w^{n},
  \qquad |w|<a .
\end{equation}
The linear term is $\psi(a)\,w$. It cannot be written as $-\zeta(1,a)\,w$, because the
Hurwitz zeta function has a pole at $s=1$. The odd terms cancel in the symmetric
difference, and the even part is exactly \eqref{eq:main1}; equivalently,
\begin{equation}\label{eq:classical}
  \log\Gamma(a+w)+\log\Gamma(a-w)-2\log\Gamma(a)
   =\sum_{k\ge1}\frac{\zeta(2k,a)}{k}\,w^{2k},\qquad |w|<a ,
\end{equation}
and it is in this form, as a sum of a series of Hurwitz zeta values, that the identity is
recorded in the literature (see, e.g., \cite{Srivastava88}) and in monographs on zeta
series \cite{SrivastavaChoi01}; the classical relation between the gamma function and the
Hurwitz zeta function goes back at least to \cite{Berndt85}. We claim no novelty
for \eqref{eq:classical}, nor for its specialization $a=1+v$, $w=u$, which is
\eqref{eq:main2}. What is new here is
(i)~two-sided remainder estimates for this series, with geometric rate $\rho$ optimal in
the sense $\frac1m\log R_m\to2\log\rho$, together with the comparison of $\rho$ and $Q$;
(ii)~the observation that the specialization $a=1+v$ converges for every $x,y>0$, which
settles Open Problems 1 and 2 of \cite{Wisniewska25} by the classical expansion;
(iii)~the localization theory of the trigamma mean $t$: the bracket \eqref{eq:locbnd},
the expansion \eqref{eq:locasym} as $c\to0$, and the bound $\theta>\frac23$ of
Theorem \ref{thm:sharp}, proved by the four-step argument of Section~\ref{sec:proofs};
(iv)~the leading large-argument asymptotics of Theorem \ref{thm:asym};
(v)~the multivariate extension of Theorem \ref{thm:multi}; and
(vi)~the applications of Section \ref{sec:app}.

Gurland's ratio has also been studied intensively by asymptotic methods. Merkle
\cite{Merkle05} derived bounds and asymptotic expansions; Alzer \cite{Alzer97} and Qi
\cite{Qi10} gave further bounds and a survey; Tian and Yang \cite{TianYang21} established
the asymptotic expansion
\[
  \frac{\Gamma(x+p)\,\Gamma(x+q)}{\Gamma\big(x+\frac{p+q}{2}\big)^{2}}
   =\exp\Big[\sum_{k=1}^{n}\frac{B_{2k}\big(\frac{1-|p-q|}{2}\big)
   -B_{2k}\big(\frac12\big)}
   {k(2k-1)\,\big(x+\frac{p+q-1}{2}\big)^{2k-1}}+\mathcal{T}_{n}\Big],
   \qquad x\to\infty ,
\]
where $p,q\in\R$ with $|p-q|\ne0$, $B_{2k}$ are the Bernoulli polynomials, and
$\mathcal{T}_{n}$ is the remainder in that expansion. The signed remainder
$(-1)^{n}\mathcal{T}_{n}$ is completely monotone when $|p-q|<1$, which yields sharp
bounds; the remainder was analysed
further by Yang and Tian \cite{YangTian24} for the
generalized ratio $\frac{\Gamma(x+a)\Gamma(x+b)}{\Gamma(x+c)\Gamma(x+d)}$ with
$a+b=c+d$. These expansions are \emph{asymptotic in $1/x$} with Bernoulli-polynomial
coefficients, whereas \eqref{eq:main1}--\eqref{eq:main2} are \emph{exact and convergent}
in powers of $\big(\frac{x-y}{2}\big)^{2}$ with Hurwitz-zeta coefficients. The two
descriptions are complementary: the former is adapted to large arguments, the latter to
arguments that are close together, and only the latter can be summed to the exact value.
Refinements of Gurland's formula for $\pi$ were given by Lin \cite{Lin13}.

Two further strands of the literature are close to the localization problem studied here.
Yang and Zheng \cite{YangZheng19} gave necessary and sufficient conditions for four classes
of ratios of gamma functions to be logarithmically completely monotonic; their
Corollary~11 and Remark~5 yield the two-sided bound \eqref{eq:YZ}, and their
Propositions~3 and~4 compare double integrals of $\psi^{(n)}$ over rectangles, which is the
same weighted-average picture that underlies \eqref{eq:locbnd}. In \cite{YangZheng16} the
mean $\psi_{n}^{-1}\big(\frac{1}{b-a}\int_{a}^{b}\psi_{n}(x+\tau)\,d\tau\big)-x$, generated
by a \emph{uniform} average of a polygamma function, is shown to increase from $\min(a,b)$
onto $\frac{a+b}{2}$; the parameter $t$ of \eqref{eq:Wt} is the corresponding mean for a
\emph{triangular} average of $\psi'$ (see Remark \ref{rem:equiv}), so the two objects are of
the same type but not the same, and neither the existence of an interior sharp constant nor
its value follows from \cite{YangZheng16}. Very recently, Moustafa and Mahmoud
\cite{MoustafaMahmoud26} obtained bounds for Gurland-type ratios in terms of the generalized
Nielsen beta function. The remainder theory of \eqref{eq:main1} and the quantitative
localization theory developed below do not appear in these works.

\subsection{Results of the present paper}
Our starting point is that \eqref{eq:jensen} is a second-order symmetric difference of
$\log\Gamma$, so the subject is governed by the Taylor expansion of $\log\Gamma$ about the
midpoint $v$. Since the singularities of $\log\Gamma$ nearest to $v>0$ are the poles of
$\Gamma$ at $0,-1,-2,\dots$, the Taylor series about $v$ converges on the disc $|z-v|<v$,
and $x=v+u$, $y=v-u$ both lie in that disc precisely when $|u|<v$, that is, precisely when
$x,y>0$. The expansion therefore converges on the whole quadrant.

\begin{theorem}[Exact expansion]\label{thm:main}
For all $x,y>0$,
\begin{align}
  \log\Gg(x,y)&=\sum_{k=1}^{\infty}\frac{1}{k}\Big(\frac{x-y}{2}\Big)^{2k}
      \zeta\Big(2k,\frac{x+y}{2}\Big),\label{eq:main1}\\[1mm]
  \log\Gs(x,y)&=\sum_{k=1}^{\infty}\frac{1}{k}\Big(\frac{x-y}{2}\Big)^{2k}
      \zeta\Big(2k,1+\frac{x+y}{2}\Big).\label{eq:main2}
\end{align}
Both series converge absolutely. Moreover \eqref{eq:main1} converges for exactly those
$(x,y)\in\R^{2}$ with $|x-y|<x+y$, and \eqref{eq:main2} for exactly those with
$|x-y|<2+x+y$; on the quadrant $x,y>0$ both hold identically, and the two regions are
optimal.
\end{theorem}

\begin{corollary}\label{cor:open12}
For every $x,y>0$ and every $m\ge1$,
\begin{align*}
  R_{m}(x,y)&=\log\Gs(x,y)-\sum_{k=1}^{m-1}\frac{1}{k}\Big(\frac{x-y}{2}\Big)^{2k}
     \zeta\Big(2k,1+\frac{x+y}{2}\Big)\\
  &=\sum_{k=m}^{\infty}\frac{1}{k}\Big(\frac{x-y}{2}\Big)^{2k}
     \zeta\Big(2k,1+\frac{x+y}{2}\Big)\;\ge\;0 .
\end{align*}
Consequently $R_{m}(x,y)\downarrow0$ and $\log\Gs(x,y)=S_{\infty}(x,y)$ for every
$x,y>0$. Open Problems 1 and 2 both have affirmative answers on the whole quadrant.
\end{corollary}

The next result quantifies the convergence. The relevant rate is
\begin{equation}\label{eq:rho}
  \rho(x,y):=\frac{|x-y|}{2+x+y}\quad(\Gs),\qquad
  \rho_{\Gg}(x,y):=\frac{|x-y|}{x+y}\quad(\Gg).
\end{equation}

\begin{theorem}[Remainder estimates]\label{thm:rem}
Let $x,y>0$, $m\ge1$, $u=\tfrac{x-y}{2}$, $v=\tfrac{x+y}{2}$,
$\rho=\rho(x,y)=|x-y|/(2+x+y)$ and
$R_m:=\log\Gs(x,y)-\sum_{k=1}^{m-1}\frac{u^{2k}}{k}\zeta(2k,1+v)$. Then
\begin{equation}\label{eq:rembnd}
  \frac{u^{2m}}{m}\,\zeta(2m,1+v)\;\le\;R_m\;\le\;
  \frac{\rho^{2m}}{m\,(1-\rho^{2})}\Big(1+\frac{1+v}{2m-1}\Big).
\end{equation}
Moreover, if $x\ne y$, then $\lim_{m\to\infty}\frac{1}{m}\log R_m=2\log\rho$, so the
geometric rate $\rho$ is optimal. The same holds for $\Gg$ with $\rho$ replaced by
$\rho_{\Gg}=|x-y|/(x+y)$ and $\zeta(2m,1+v)$ by $\zeta(2m,v)$.
\end{theorem}

The limit fixes the base of the exponential. It does not assert that the prefactors in
\eqref{eq:rembnd} are best possible: the last row of Table \ref{tab:rem} has lower bound
$0.3147$ and upper bound $15.9456$.

\begin{proposition}\label{prop:compare}
For all $x,y>0$ one has $\rho(x,y)\le Q(x,y)$, with equality iff $x=y$. Moreover
$\rho(x,y)<1$ for all $x,y>0$, whereas $\sup_{x,y>0}Q(x,y)=+\infty$; for instance
$Q(x,1)=\frac{x-1}{2(1+\sqrt x)}\sim\frac{\sqrt x}{2}$ as $x\to\infty$. Hence $\{Q<1\}$ is a
proper subset of the quadrant on which \eqref{eq:main2} converges.
\end{proposition}

\begin{corollary}\label{cor:bil}
For all $x,y>0$ and all $m\ge1$,
\[
  S_m(x,y)\;\le\;\log\Gs(x,y)\;\le\;S_m(x,y)+
  \frac{\rho^{2m}}{m(1-\rho^{2})}\Big(1+\frac{1+v}{2m-1}\Big),
\]
where $S_m(x,y)=\sum_{k=1}^{m-1}\frac{1}{k}\big(\frac{x-y}{2}\big)^{2k}\zeta(2k,1+v)$.
Taking $m=2$ improves the upper bound of \eqref{eq:Wbil} whenever
$\rho<\frac{\sqrt5}{3}$, while $S_m$ increases to $\log\Gs(x,y)$ as $m\to\infty$.
\end{corollary}

We also determine the behaviour of the ratio at large arguments. Recall
\begin{equation}\label{eq:Lambda}
  \Lambda(\rho):=(1+\rho)\log(1+\rho)+(1-\rho)\log(1-\rho),\qquad |\rho|<1 .
\end{equation}
One has $\Lambda(\rho)=\sum_{k\ge1}\rho^{2k}/\big(k(2k-1)\big)$, $\Lambda(\rho)>0$ for
$\rho\ne0$, and $\Lambda(\rho)=2\rho\artanh\rho+\log(1-\rho^{2})$. Let $H_b$ denote the
binary entropy
\begin{equation}\label{eq:entropy}
  H_b(p)=-p\log p-(1-p)\log(1-p),
\end{equation}
formed with the same logarithm as above. For $p=\frac{1+\rho}{2}$,
\begin{align*}
  H_b(p)
  &=-\frac{1+\rho}{2}\log\frac{1+\rho}{2}-\frac{1-\rho}{2}\log\frac{1-\rho}{2}
   =\log 2-\frac12\Lambda(\rho),
\end{align*}
and therefore
\begin{equation}\label{eq:entropyid}
  \Lambda(\rho)=2\Big(\log 2-H_b\Big(\frac{1+\rho}{2}\Big)\Big).
\end{equation}
As $\Lambda$ vanishes to second order at the origin it is convenient to divide it by
$\rho^{2}$. Set
\begin{equation}\label{eq:ell}
  \ell(\rho):=\frac{\Lambda(\rho)}{\rho^{2}}
   =\sum_{k\ge1}\frac{\rho^{2k-2}}{k(2k-1)}\quad(0<\rho<1),
  \qquad \ell(0):=1 .
\end{equation}
The series converges for $|\rho|<1$, so $\ell$ is continuous and strictly increasing on
$[0,1)$, with values at least $1$, and $\Lambda(\rho)=\rho^{2}\ell(\rho)$. The binary
entropy enters the asymptotics through \eqref{eq:entropyid}; the coefficient in
Theorem \ref{thm:asym} is the normalization $\ell$.

\begin{theorem}[Large-argument asymptotics]\label{thm:asym}
Let $x_{j},y_{j}>0$, $v_{j}=\tfrac{x_j+y_j}{2}\to\infty$, $u_{j}=\tfrac{x_j-y_j}{2}$, and
put $\rho_{j}:=\frac{|u_{j}|}{1+v_{j}}$. Suppose
\[
  \rho_{j}\;\longrightarrow\;\rho_{0}\in[0,1).
\]
Then, with $\ell$ as in \eqref{eq:ell},
\[
  \log\Gs(x_{j},y_{j})
   =\frac{u_{j}^{2}}{1+v_{j}}\,\ell(\rho_{j})\big(1+o(1)\big)
   =(1+v_{j})\,\Lambda(\rho_{j})\big(1+o(1)\big)
  \qquad (j\to\infty),
\]
and the analogous statement holds for $\log\Gg$ with $1+v_j$ replaced by $v_j$
throughout, i.e.\ under the hypothesis $|u_j|/v_j\to\rho_0\in[0,1)$ one has
$\log\Gg=\frac{u_j^{2}}{v_j}\ell\big(\frac{|u_j|}{v_j}\big)(1+o(1))$. (The sign of $u_j$ is
immaterial, since $\Lambda$ and $\ell$ are even.)
\end{theorem}

The expansion extends verbatim to any number of variables.

\begin{theorem}[Multivariate expansion]\label{thm:multi}
Let $n\ge2$, $x_{1},\dots,x_{n}>0$ and $\bar x=\frac1n\sum_{i=1}^{n}x_{i}$. Write
$\mu_{j}=\sum_{i=1}^{n}(x_{i}-\bar x)^{j}$ for $j\ge2$. If
$\max_{i}|x_{i}-\bar x|<\bar x$, then
\begin{equation}\label{eq:multi}
  \log\frac{\prod_{i=1}^{n}\Gamma(x_{i})}{\Gamma(\bar x)^{n}}
   =\sum_{j=2}^{\infty}\frac{(-1)^{j}}{j}\,\zeta(j,\bar x)\,\mu_{j},
\end{equation}
the series converging absolutely. In particular
$\prod_{i=1}^{n}\Gamma(x_{i})\ge\Gamma(\bar x)^{n}$, with equality iff $x_{1}=\dots=x_{n}$.
\end{theorem}

Finally we address Open Problem 3. The parameter $t$ is unique, the gap $v-t$ admits an
explicit bracket and an expansion in powers of $c^{2}$, and the ratio $\theta$ is bounded
by the constant $\frac23$.

\begin{theorem}[Localization of $t$]\label{thm:loc}
Let $x\ne y$ be positive, $v=\tfrac{x+y}{2}$, $s=\sqrt{xy}$, $c=\tfrac{|x-y|}{2}$,
$A=1+v$ and $\rho=c/A$. Then:
\begin{enumerate}
\item[\rm(i)] there is a unique $t=t(x,y)\in(s,v)$ satisfying \eqref{eq:Wt};
\item[\rm(ii)] with $\rho=\rho(x,y)$ one has
  \begin{equation}\label{eq:locbnd}
    \frac{c^{2}\,\zeta(4,A)}{4\,\zeta(3,A)+6v\,\zeta(4,A)}
    \;\le\; v-t \;\le\;
    \frac{c^{2}\,\zeta(4,A)}{4\,\zeta(3,A)}\Big(1+\frac{2\rho^{2}}{3(1-\rho^{2})}\Big);
  \end{equation}
\item[\rm(iii)] as $c\to0^{+}$ with $v$ fixed,
  \begin{equation}\label{eq:locasym}
    v-t=\frac{c^{2}\zeta(4,A)}{4\zeta(3,A)}
     +\frac{c^{4}}{2\zeta(3,A)}\Big[\frac{\zeta(6,A)}{3}
       -\frac{3\,\zeta(4,A)^{3}}{16\,\zeta(3,A)^{2}}\Big]+O(c^{6}).
  \end{equation}
\end{enumerate}
\end{theorem}

The coefficients in \eqref{eq:locbnd} give the scale of $v-t$. They are not optimal: at
$(x,y)=(0.1,199.9)$ one has $v-t=26.6699$, while the upper bound equals $514.83$
(Section \ref{sec:num}). The position of $t$ between $s$ and $v$ is settled by the next
theorem.

\begin{theorem}[Sharp localization]\label{thm:sharp}
For all $x\ne y>0$ the ratio \eqref{eq:theta} satisfies
\begin{equation}\label{eq:conj}
  \theta(x,y)\;>\;\frac{2}{3},
\end{equation}
equivalently $t(x,y)>\frac13\big(x+y+\sqrt{xy}\big)$. The constant $\frac23$ is best
possible and the bound is not attained: $\theta(x,y)$ tends to $\frac23$
along $v\to\infty$ with $|x-y|/(x+y)\to0$.
\end{theorem}

\begin{remark}[The midpoint inequality]\label{rem:mid}
Since $\frac23>\frac12$, Theorem \ref{thm:sharp} gives
\begin{equation}\label{eq:midpoint}
  t(x,y)>\frac12\Big(\sqrt{xy}+\frac{x+y}{2}\Big)
\end{equation}
for every $x\ne y>0$. The same bound also follows from the upper estimate in
\eqref{eq:locbnd}, on a subset of the quadrant, by the following comparison. That upper
bound is strictly smaller than $\frac{v-s}{2}=\frac{c^{2}}{2(v+s)}$ as soon as
\[
  \frac{(v+s)\,\zeta(4,A)}{2\,\zeta(3,A)}
  \Big(1+\frac{2\rho^{2}}{3(1-\rho^{2})}\Big)<1 .
\]
The inequality $(n+A)^{-1}\le A^{-1}$ gives $\zeta(4,A)\le A^{-1}\zeta(3,A)$, and hence
$(v+s)\zeta(4,A)\le\frac{v+s}{A}\zeta(3,A)$. The display therefore holds whenever
\begin{equation}\label{eq:condA}
  \frac{v+s}{2(1+v)}\Big(1+\frac{2\rho^{2}}{3(1-\rho^{2})}\Big)<1 .
\end{equation}
A second sufficient condition uses the summands of $\Phi:=\log\Gs(x,y)/c^{2}$. Let
$w=\frac{v-s}{2}$ and
$K(\xi)=\frac{1}{c^{2}}\log\frac{1}{1-c^{2}/\xi^{2}}$ for $\xi>c$, and write
$M(\xi)=(\xi-w)^{-2}$. Then $\sum_{n\ge0}K(A+n)=\Phi$ and
$\sum_{n\ge0}M(A+n)=\zeta\big(2,1+\frac{v+s}{2}\big)$, so
$t>\frac{v+s}{2}$ if and only if $\sum_{n\ge0}(K-M)(A+n)<0$. On
$\big[v+\frac12,\infty\big)$ one has $\xi>c$, and both functions are positive, decreasing
and convex: $M''(\xi)=6(\xi-w)^{-4}>0$, while
$K'(\xi)=-\frac{2}{\xi^{3}-c^{2}\xi}$ and
$K''(\xi)=\frac{2(3\xi^{2}-c^{2})}{(\xi^{3}-c^{2}\xi)^{2}}>0$ since $\xi>c$. For a convex decreasing $f$
with $f(\infty)=0$, the midpoint rule gives
$\sum_{n\ge0}f(A+n)\le\int_{A-1/2}^{\infty}f$ and the trapezoid rule gives
$\sum_{n\ge0}f(A+n)\ge\int_{A}^{\infty}f+\frac{f(A)}{2}$. Applying the first to $K$ and
the second to $M$,
\[
  \sum_{n\ge0}(K-M)(A+n)
  \;\le\;\int_{v+1/2}^{\infty}K-\frac{1}{A-w}-\frac{1}{2(A-w)^{2}},
\]
because $\int_{A}^{\infty}M=\frac{1}{A-w}$. Consequently $t>\frac{v+s}{2}$ whenever
\begin{equation}\label{eq:condB}
  \int_{v+1/2}^{\infty}K-\frac{1}{A-w}-\frac{1}{2(A-w)^{2}}<0 .
\end{equation}
The integral equals
\begin{equation}\label{eq:intK}
  \int_{B}^{\infty}K
  =\frac{1}{c^{2}}\Big[-B\log\frac{B^{2}}{B^{2}-c^{2}}
     +c\log\frac{B+c}{B-c}\Big],\qquad B>c .
\end{equation}
The proof of Theorem \ref{thm:sharp} is independent of \eqref{eq:condA} and of
\eqref{eq:condB}.
\end{remark}

\begin{remark}[A product formula]\label{rem:product}
The Weierstrass product for $\Gamma$ gives a representation of $\Gs$ that makes the
identity $\Phi=\sum_{m\ge0}K(A+m)$ transparent, where
$K(\xi)=\frac{1}{c^{2}}\log\frac{1}{1-c^{2}/\xi^{2}}$ and $\Phi=\log\Gs/c^{2}$. Indeed
$\Gamma(z)=\frac{e^{-\gamma z}}{z}\prod_{n\ge1}\frac{e^{z/n}}{1+z/n}$ yields, with
$v=\frac{x+y}{2}$ and $c=\frac{|x-y|}{2}$,
\[
  \frac{\Gamma(v+c)\Gamma(v-c)}{\Gamma(v)^{2}}
   =\frac{v^{2}}{v^{2}-c^{2}}\prod_{n\ge1}\frac{(n+v)^{2}}{(n+v)^{2}-c^{2}} ,
\]
the factors $e^{\pm z/n}$ cancelling. Since
$\Gs(x,y)=\frac{4xy}{(x+y)^{2}}\Gg(x,y)=\frac{v^{2}-c^{2}}{v^{2}}\cdot
 \frac{\Gamma(v+c)\Gamma(v-c)}{\Gamma(v)^{2}}$, the factors $v^{2}-c^{2}$ cancel as well
and
\begin{equation}\label{eq:product}
  \Gs(x,y)=\prod_{n\ge1}\frac{(n+v)^{2}}{(n+v)^{2}-c^{2}}
          =\prod_{m\ge0}\frac{(A+m)^{2}}{(A+m)^{2}-c^{2}} ,
\end{equation}
with $A=1+v$. Taking logarithms recovers
\[
  \log\Gs=\sum_{m\ge0}\log\frac{(A+m)^{2}}{(A+m)^{2}-c^{2}}
  =c^{2}\sum_{m\ge0}K(A+m)
  =\sum_{k\ge1}\frac{c^{2k}}{k}\zeta(2k,A),
\]
which is Theorem~\ref{thm:main} for $\Gs$. Each factor of \eqref{eq:product} exceeds $1$,
because $c<v<A$, and this is the Jensen--Gurland inequality in multiplicative form.
\end{remark}

\begin{remark}[Scaling limit]\label{rem:scaling}
Fix $\rho=\frac{c}{A}$ and let $A\to\infty$. Since $\zeta(2m+2,A)\sim\frac{A^{-(2m+1)}}{2m+1}$
and $v-s=A\big(1-\sqrt{1-\rho^{2}}\big)(1+o(1))$, the comparison
$\zeta(2,A-\Delta)>\Phi$ behind Theorem \ref{thm:sharp} becomes, at leading order in
$A^{-1}$,
\begin{equation}\label{eq:scaling}
  \frac{\kappa}{1-\kappa}\;>\;P(\rho^{2})-1,
  \qquad \kappa=\frac{1-\sqrt{1-\rho^{2}}}{3},
\end{equation}
where $P$ is the series of Step~2 in the proof of Theorem \ref{thm:sharp},
\begin{equation}\label{eq:P}
  P(y)=\sum_{m\ge0}\frac{y^{m}}{(m+1)(2m+1)}\qquad(0\le y<1),
\end{equation}
with the closed form
$P(\rho^{2})=2\frac{\artanh\rho}{\rho}+\frac{\log(1-\rho^{2})}{\rho^{2}}$. The difference
of the two sides of \eqref{eq:scaling},
$D(\rho)=\frac{\kappa}{1-\kappa}-P(\rho^{2})+1$, is positive on $(0,1)$, and this is a
consequence of Step~2 of that proof. Indeed, with $y=\rho^{2}$,
\[
  \frac{\kappa}{1-\kappa}=\frac{1-\sqrt{1-y}}{2+\sqrt{1-y}}=h(y),
  \qquad P(\rho^{2})=P(y),
  \qquad\text{so } D=h-P+1 .
\]
Step~2 writes $h=\sum_{m\ge1}a_{m}y^{m}$ and $P=1+\sum_{m\ge1}b_{m}y^{m}$ with
$a_{1}=b_{1}=\frac16$ and $a_{m}>b_{m}$ for every $m\ge2$; hence
$D=\sum_{m\ge2}(a_{m}-b_{m})y^{m}$ has positive coefficients, which proves $D>0$ and at the
same time identifies the expansion
$D(\rho)=\frac{\rho^{4}}{360}+\frac{11\rho^{6}}{3024}+O(\rho^{8})$, since
$a_2-b_2=\frac1{360}$ and $a_3-b_3=\frac{11}{3024}$ are the coefficients computed there
(with a positive $O(\rho^{8})$ tail). Numerically $D(\rho)\approx\rho^{4}/360$ down to
$\rho=5\cdot10^{-4}$ (where $D=1.74\cdot10^{-16}$) and $D(0.99)=5.10\cdot10^{-2}$. The
constant $\frac23$ is decided by an $O(\rho^{4})$ cancellation, which is why a first-order
estimate does not reach Theorem~\ref{thm:sharp}.
\end{remark}

\begin{remark}\label{rem:equiv}
Theorem \ref{thm:sharp} can be restated as a comparison of $\zeta(2,\cdot)$ with its
own triangular average. With $A=1+v$, $c=|x-y|/2$,
$K(\xi)=\frac{1}{c^{2}}\log\frac{1}{1-c^{2}/\xi^{2}}$ and
$\Phi=\log\Gs(x,y)/c^{2}$, the identity $\sum_{n\ge0}K(A+n)=\Phi$ of
Remark \ref{rem:product} combines with
$\psi(A+b)-\psi(A-b)=\int_{-b}^{b}\psi'(A+z)\,dz$ to give
\[
  \Phi=\frac{1}{c^{2}}\int_{0}^{c}\big[\psi(A+b)-\psi(A-b)\big]\,db
     =\int_{-c}^{c}\psi'(A+z)\,\frac{c-|z|}{c^{2}}\,dz .
\]
Since $\zeta(2,1+t)=\psi'(1+t)$, the inequality of Theorem \ref{thm:sharp} is therefore
equivalent to
\[
  \int_{-c}^{c}\Big[\psi'\!\big(A-\Delta\big)-\psi'(A+z)\Big]\frac{c-|z|}{c^{2}}\,dz>0,
  \qquad \Delta=\frac{v-s}{3}=\frac{c^{2}}{3(v+s)},
\]
where $s=\sqrt{xy}$: it compares the value of $\psi'=\zeta(2,\cdot)$ at the point
$A-\Delta$ with its average over $[-c,c]$ against the triangular weight $(c-|z|)/c^{2}$,
whose variance is $c^{2}/6$. Expanding both sides about $A$, the terms of order $c^{0}$ and
$c^{1}$ cancel (the weight is even), and the first surviving comparison is
\[
  \frac{\zeta(4,A)}{\zeta(3,A)}<\frac{4}{3(v+s)} ,
\]
which holds for all $x,y>0$, with relative margin $\frac{1}{A}$ because
$\zeta(4,A)/\zeta(3,A)=\frac{2}{3A}\big(1+O(A^{-1})\big)$ and $s\le v$. The decision rests
on the next terms, of relative size $\rho^{2}$, and this is why the constant $\frac23$ is
sharp and is not attained.
\end{remark}

\begin{remark}\label{rem:limit}
The limiting statement in Theorem \ref{thm:sharp} follows from Theorem
\ref{thm:loc}(iii). Take $x_{j}=v_{j}(1+\eps_{j})$, $y_{j}=v_{j}(1-\eps_{j})$ with
$v_{j}\to\infty$, $\eps_{j}\to0$ and $v_{j}\eps_{j}\to0$. Then $c_{j}=v_{j}\eps_{j}\to0$,
so \eqref{eq:locasym} applies and, since $\zeta(4,A)/\zeta(3,A)\sim2/(3A)$ as
$A\to\infty$,
\[
  v_{j}-t_{j}\sim\frac{v_{j}^{2}\eps_{j}^{2}}{4}\cdot\frac{2}{3(1+v_{j})}
  \sim\frac{v_{j}\eps_{j}^{2}}{6},
  \qquad
  v_{j}-s_{j}=\frac{v_{j}\eps_{j}^{2}}{1+\sqrt{1-\eps_{j}^{2}}}\sim\frac{v_{j}\eps_{j}^{2}}{2}.
\]
Hence $(v_{j}-t_{j})/(v_{j}-s_{j})\to\frac13$ and
$(t_{j}-s_{j})/(v_{j}-s_{j})\to\frac23$. Along this path
$\theta=\frac23+\frac{1}{6v}+O(v^{-2})$. The column $\eps=v^{-2}$ of Table \ref{tab:inf} is such a
path and confirms the expansion to nine significant digits. A fixed $\eps$ produces a
larger limit: $\theta\to\frac23+g(\eps)$ with $g(\eps)>0$. Numerically
$g(\eps)\sim\eps^{2}/180$ as $\eps\to0$; for instance $g(10^{-1})=5.60\cdot10^{-5}$, while
$10^{-2}/180=5.56\cdot10^{-5}$.
\end{remark}

\section{Proofs}\label{sec:proofs}

\subsection{Proof of Theorem \ref{thm:main}}
Let $v>0$ and put $\Omega=\{z\in\mathbb{C}:\operatorname{Re}z>0\}$. The principal branch of
$\log\Gamma$ is holomorphic on $\Omega$, and the disc $D(v,v)=\{z:|z-v|<v\}$ is contained in
$\Omega$. The nearest singularities of $\log\Gamma$ to $v$ are the poles of $\Gamma$ at
$0,-1,-2,\dots$, all at distance $\ge v$ from $v$, and $|0-v|=v$; hence the Taylor series of
$\log\Gamma$ about $v$ has radius of convergence exactly $v$ and converges to $\log\Gamma$
on $D(v,v)$.

Fix $u$ with $|u|<v$. Then $v\pm u\in D(v,v)$ and
\[
  \log\Gamma(v+w)=\log\Gamma(v)+\sum_{n=1}^{\infty}\frac{\psi^{(n-1)}(v)}{n!}\,w^{n},
  \qquad |w|<v ,
\]
where $\psi=\Gamma'/\Gamma$ and $\psi^{(n)}$ are the polygamma functions. Evaluating at
$w=u$ and $w=-u$ and adding,
\[
  \log\Gamma(v+u)+\log\Gamma(v-u)-2\log\Gamma(v)
   =2\sum_{k=1}^{\infty}\frac{\psi^{(2k-1)}(v)}{(2k)!}\,u^{2k}.
\]
The relation $\psi^{(n)}(z)=(-1)^{n+1}n!\,\zeta(n+1,z)$ for $\operatorname{Re}z>0$
(see e.g.\ \cite[\S5.15]{DLMF}) gives, for $n=2k-1$,
$\psi^{(2k-1)}(v)=(2k-1)!\,\zeta(2k,v)$. Therefore
\[
  \log\Gg(x,y)=2\sum_{k=1}^{\infty}\frac{(2k-1)!}{(2k)!}\,\zeta(2k,v)\,u^{2k}
   =\sum_{k=1}^{\infty}\frac{u^{2k}}{k}\,\zeta(2k,v),
\]
which is \eqref{eq:main1}; its region of convergence is $|u|<v$, i.e.\
$\{|x-y|<x+y\}$, and since all terms are non-negative the convergence is absolute. The
same argument applied to $z\mapsto\log\Gamma(1+z)$ at the point $1+v$, whose singularities
are $\{-1,-2,-3,\dots\}$ and hence lie at distance $\ge1+v$ from $1+v$, with equality at
$-1$, yields \eqref{eq:main2} with convergence exactly for $|u|<1+v$, i.e.\ for
$|x-y|<2+x+y$. Finally, $|x-y|<x+y$ and $|x-y|<2+x+y$ both hold for all $x,y>0$; the first
fails for $(x,y)=(1,-\tfrac12)$ and the second for $(x,y)=(3,-2)$, so the stated regions are
optimal. \qed

\subsection{Proof of Corollary \ref{cor:open12}}
By Theorem \ref{thm:main}, for every $x,y>0$ the series \eqref{eq:main2} converges and all
its terms are positive. Hence
$R_{m}(x,y)=\sum_{k=m}^{\infty}\frac{u^{2k}}{k}\zeta(2k,1+v)$
is the tail of a convergent series of positive terms; it is positive and decreases to $0$.
This is Open Problems 1 and 2. \qed

\subsection{Proof of Theorem \ref{thm:rem}}
Put $A=1+v$, $c=|u|$, $\rho=c/A\in(0,1)$. We use the integral-test estimate
\begin{equation}\label{eq:zetabnd}
  \zeta(2k,A)\;\le\;A^{-2k}\Big(1+\frac{A}{2k-1}\Big),\qquad k\ge1,\ A>0,
\end{equation}
obtained from
$\zeta(2k,A)=\sum_{n\ge0}(n+A)^{-2k}\le A^{-2k}+\int_{0}^{\infty}(\tau+A)^{-2k}\,d\tau$.

\emph{Lower bound.} All terms of the tail are positive, so
$R_m\ge\frac{u^{2m}}{m}\zeta(2m,A)$.

\emph{Upper bound.} By \eqref{eq:zetabnd},
\[
  R_m\le\sum_{k=m}^{\infty}\frac{c^{2k}}{k}A^{-2k}\Big(1+\frac{A}{2k-1}\Big)
   \le\Big(1+\frac{A}{2m-1}\Big)\sum_{k=m}^{\infty}\frac{\rho^{2k}}{k}
   =\Big(1+\frac{A}{2m-1}\Big)\rho^{2m}\sum_{j=0}^{\infty}\frac{\rho^{2j}}{m+j},
\]
and $\sum_{j\ge0}\frac{\rho^{2j}}{m+j}\le\frac{1}{m(1-\rho^{2})}$, which is
\eqref{eq:rembnd} for $\Gs$.

\emph{Optimality.} Let $x\ne y$, so $c>0$ and $\rho\in(0,1)$. Write
$\zeta(2m,A)=A^{-2m}(1+\omega_m)$; since
$\sum_{n\ge1}\big(\frac{A}{n+A}\big)^{2m}
\le\int_0^{\infty}\big(\frac{A}{A+\tau}\big)^{2m}\,d\tau
=\frac{A}{2m-1}$, we have $0\le\omega_m\le\frac{A}{2m-1}\to0$. Hence
$R_m\ge\frac{c^{2m}}{m}\zeta(2m,A)=\frac{1+\omega_m}{m}\rho^{2m}$, so
$\frac1m\log R_m\ge2\log\rho+o(1)$; conversely \eqref{eq:rembnd} gives
$\frac1m\log R_m\le2\log\rho+\frac1m\log\frac{1+\frac{A}{2m-1}}{m(1-\rho^{2})}
=2\log\rho+o(1)$. The proof for $\Gg$ is identical with $A$ replaced by $v$; there
$\zeta(2k,v)\le v^{-2k}\big(1+\frac{v}{2k-1}\big)$ and $\rho_{\Gg}=c/v$. \qed

\subsection{Proof of Proposition \ref{prop:compare}}
Since $x+y\ge2\sqrt{xy}$, $\rho=\frac{|x-y|}{2+x+y}\le\frac{|x-y|}{2+2\sqrt{xy}}=Q$, with
equality iff $x=y$. Also $|x-y|<x+y<2+x+y$, so $\rho<1$. Finally
$Q(x,1)=\frac{x-1}{2(1+\sqrt x)}\sim\frac{\sqrt x}{2}\to\infty$. \qed

\subsection{Proof of Corollary \ref{cor:bil}}
The two-sided inequality is Theorem \ref{thm:rem} restated, and
$S_{m}\to\log\Gs(x,y)$ as $m\to\infty$ by Corollary \ref{cor:open12}. It remains to
compare, at $m=2$, the correction in \eqref{eq:rembnd} with the gain over the upper
bound of \eqref{eq:Wbil}. Write $s=\sqrt{xy}$, $c=|u|$ and $A=1+v$, so that
$\rho=c/A$ and
\[
  c^{2}\big[\zeta(2,1+s)-\zeta(2,1+v)\big]
  =2c^{2}\int_{1+s}^{1+v}\zeta(3,z)\,dz
  \;\ge\;2c^{2}\int_{1+s}^{1+v}\frac{dz}{2z^{2}}
  =\frac{c^{2}(v-s)}{(1+s)(1+v)},
\]
because $\zeta(3,z)=\sum_{n\ge0}(n+z)^{-3}\ge\int_{0}^{\infty}(\tau+z)^{-3}\,d\tau=\frac{1}{2z^{2}}$
(the integrand decreases, so $\int_{n}^{n+1}(\tau+z)^{-3}\,d\tau\le(n+z)^{-3}$ for every $n\ge0$).
Since $v-s=\frac{v^{2}-s^{2}}{v+s}=\frac{c^{2}}{v+s}$ and $(v+s)(1+s)\le2v(1+v)$, the gain
is at least $\frac{c^{4}}{2v(1+v)^{2}}$. The correction in \eqref{eq:rembnd} at $m=2$ is
\[
  \frac{\rho^{4}}{2(1-\rho^{2})}\Big(1+\frac{A}{3}\Big)
  =\frac{c^{4}(A+3)}{6A^{4}(1-\rho^{2})},
\]
which is smaller as soon as $\frac{v(A+3)}{3A^{2}(1-\rho^{2})}<1$. As
$\frac{v(A+3)}{3A^{2}}=\frac{v(v+4)}{3(v+1)^{2}}\le\frac49$, with equality at $v=2$, this
holds whenever $\rho<\frac{\sqrt5}{3}$. \qed

\subsection{Proof of Theorem \ref{thm:asym}}
If $c_{j}=0$, i.e.\ $x_{j}=y_{j}$ for some $j$, then $\log\Gs(x_{j},y_{j})=0$ and the
assertion is trivial, so we may assume $c_{j}>0$. Put $A_{j}=1+v_{j}$, $c_{j}=|u_{j}|$,
$\rho_{j}=c_{j}/A_{j}\in(0,1)$, so $\rho_{j}\to\rho_{0}$ and
$\Phi_{j}:=\frac{\log\Gs(x_{j},y_{j})}{c_{j}^{2}}
=\sum_{k\ge1}\frac{c_{j}^{2k-2}}{k}\zeta(2k,A_{j})$. Set
$\beta_{k}(A):=\frac{A^{2k-1}}{k}\zeta(2k,A)$, so that
$A_{j}\Phi_{j}=\sum_{k\ge1}\rho_{j}^{2k-2}\beta_{k}(A_{j})$. Since
$\frac{A^{2k-1}}{(n+A)^{2k}}=\frac1A\big(\frac{A}{n+A}\big)^{2k}$ and
$\tau\mapsto\big(\frac{A}{A+\tau}\big)^{2k}$ is decreasing, the integral test gives, for all
$k\ge1$ and $A>0$,
\begin{equation}\label{eq:EM}
  0\;\le\;\beta_{k}(A)\;=\;\frac{1}{kA}\sum_{n\ge0}\Big(\frac{A}{n+A}\Big)^{2k}
  \;\le\;\frac{1}{kA}\Big(1+\int_{0}^{\infty}\Big(\frac{A}{A+\tau}\Big)^{2k}\,d\tau\Big)
  \;=\;\frac{1}{kA}+\frac{1}{k(2k-1)} .
\end{equation}
Hence, for any $\rho_{1}$ with $\rho_{0}<\rho_{1}<1$ and all large $j$,
\[
  \rho_{j}^{2k-2}\beta_{k}(A_{j})\;\le\;\rho_{1}^{2k-2}\Big(\frac{1}{k}+\frac{1}{k(2k-1)}\Big)
  \qquad(k\ge1),
\]
and the dominating series converges, so dominated convergence applies. On the other hand,
for each fixed $k$ the Riemann sum $\frac1A\sum_{n\ge0}\big(\frac{A}{n+A}\big)^{2k}$
converges to $\int_{0}^{\infty}(1+\tau)^{-2k}\,d\tau=\frac{1}{2k-1}$ as $A\to\infty$, so
$\beta_{k}(A)\to\frac{1}{k(2k-1)}$. Therefore, by the definition \eqref{eq:ell} of $\ell$,
\[
  A_{j}\Phi_{j}\;\longrightarrow\;\sum_{k\ge1}\frac{\rho_{0}^{2k-2}}{k(2k-1)}
   =\ell(\rho_{0}),
\]
the case $\rho_{0}=0$ being the definition $\ell(0)=1$, and the case $\rho_{0}>0$ following
from the identity $\sum_{k\ge1}z^{k-1}/\big(k(2k-1)\big)=\Lambda(\sqrt z)/z$ for $0<z<1$,
which is obtained by integrating $\sum_{k\ge1}z^{k-1}/(2k-1)=\artanh\sqrt z/\sqrt z$ term by
term. Since $c_{j}=A_{j}\rho_{j}$, the identity $\log\Gs=c_{j}^{2}\Phi_{j}$ reads
$\log\Gs=\rho_{j}^{2}A_{j}\cdot A_{j}\Phi_{j}$; as $\ell$ is continuous at $\rho_{0}$,
\[
  \log\Gs(x_{j},y_{j})
   =\frac{u_{j}^{2}}{1+v_{j}}\,\ell(\rho_{j})\big(1+o(1)\big)
   =(1+v_{j})\,\Lambda(\rho_{j})\big(1+o(1)\big),
\]
the second equality being $\Lambda(\rho)=\rho^{2}\ell(\rho)$. The proof for $\Gg$ is the same
with $A_{j}$ replaced by $v_{j}$ and $\rho_{j}$ by $|u_{j}|/v_{j}$, and gives
$\log\Gg=\frac{u_{j}^{2}}{v_{j}}\ell\big(\frac{|u_{j}|}{v_{j}}\big)(1+o(1))$. \qed

\subsection{Proof of Theorem \ref{thm:multi}}
Each $x_{i}$ satisfies $|x_{i}-\bar x|<\bar x$, so $x_{i}$ lies in the disc
$D(\bar x,\bar x)\subset\{\operatorname{Re}z>0\}$ on which $\log\Gamma$ is holomorphic with
radius of convergence exactly $\bar x$. Hence
\[
  \log\Gamma(x_{i})=\log\Gamma(\bar x)+\sum_{m\ge1}\frac{\psi^{(m-1)}(\bar x)}{m!}
   (x_{i}-\bar x)^{m}
\]
converges for every $i$. Summing over $i$ and using $\sum_{i}(x_{i}-\bar x)=0$,
\[
  \sum_{i=1}^{n}\log\Gamma(x_{i})-n\log\Gamma(\bar x)
  =\sum_{m\ge2}\frac{\psi^{(m-1)}(\bar x)}{m!}\,\mu_{m}
  =\sum_{m\ge2}\frac{(-1)^{m}}{m}\zeta(m,\bar x)\,\mu_{m},
\]
which is \eqref{eq:multi}; absolute convergence follows from
$|\mu_{m}|\le n\max_{i}|x_{i}-\bar x|^{m}$ and $|x_{i}-\bar x|<\bar x$. The inequality is
Jensen's inequality for the convex function $\log\Gamma$, with equality only when all
$x_{i}$ coincide. \qed

\subsection{Proof of Theorem \ref{thm:loc}}

\emph{(i) Uniqueness.} The function $g(z)=\zeta(2,1+z)$ is continuous and strictly decreasing
on $(0,\infty)$. By \eqref{eq:Wbil} and the positivity of the terms of \eqref{eq:main2},
\[
  \Big(\frac{x-y}{2}\Big)^{2}\zeta(2,1+v)\le\log\Gs(x,y)\le
  \Big(\frac{x-y}{2}\Big)^{2}\zeta(2,1+s),
\]
so $\Phi:=\log\Gs(x,y)/c^{2}$ satisfies $g(v)\le\Phi\le g(s)$, both inequalities being
strict for $x\ne y$. Hence exactly one $t\in(s,v)$ satisfies $g(t)=\Phi$.

\emph{(ii) The bracket.} Write $\delta=v-t\in(0,v-s)$ and $A=1+v$. From \eqref{eq:main2},
\begin{equation}\label{eq:keyid}
  \Phi=\zeta(2,A)+\frac{c^{2}}{2}\zeta(4,A)+W,\qquad
   W:=\sum_{k\ge3}\frac{c^{2k-2}}{k}\zeta(2k,A)\ge0 .
\end{equation}
Taylor expansion of $z\mapsto\zeta(2,z)$ about $A$, all of whose terms are positive, gives
\begin{equation}\label{eq:taylor}
  \zeta(2,A-\delta)=\zeta(2,A)+2\delta\zeta(3,A)+3\delta^{2}\zeta(4,A)+U,
  \qquad U:=\sum_{j\ge3}(j+1)\delta^{j}\zeta(j+2,A)\ge0 .
\end{equation}
Since $\zeta(2,A-\delta)=\Phi$, subtraction yields
\begin{equation}\label{eq:balance}
  2\delta\zeta(3,A)+3\delta^{2}\zeta(4,A)+U=\frac{c^{2}}{2}\zeta(4,A)+W .
\end{equation}
For the lower bound write $\zeta_{j}:=\zeta(j,A)$ and put
\begin{equation}\label{eq:d0}
  \delta_{0}:=\frac{c^{2}\zeta_{4}}{4\zeta_{3}+6v\zeta_{4}} .
\end{equation}
We prove $\Phi>\zeta(2,A-\delta_{0})$, which by strict decrease of $z\mapsto\zeta(2,z)$
gives $\delta>\delta_{0}$, i.e.\ the lower bound in \eqref{eq:locbnd}. By \eqref{eq:keyid}
we have $\Phi\ge\zeta_{2}+\frac{c^{2}}{2}\zeta_{4}$ (as $W\ge0$), while by
\eqref{eq:taylor},
$\zeta(2,A-\delta_{0})=\zeta_{2}+2\delta_{0}\zeta_{3}+3\delta_{0}^{2}\zeta_{4}+U_{0}$
with $U_{0}:=\sum_{j\ge3}(j+1)\delta_{0}^{j}\zeta_{j+2}\ge0$. It therefore suffices to
verify
\begin{equation}\label{eq:lbmargin}
  U_{0}\;<\;\frac{c^{2}}{2}\zeta_{4}-2\delta_{0}\zeta_{3}-3\delta_{0}^{2}\zeta_{4}
  \;=\;3\,\delta_{0}\zeta_{4}\,(v-\delta_{0}),
\end{equation}
the equality following from \eqref{eq:d0} together with
$\delta_{0}(2\zeta_{3}+3v\zeta_{4})=\frac{c^{2}}{2}\zeta_{4}$.

Now $\delta_{0}<\frac{c^{2}\zeta_{4}}{6v\zeta_{4}}=\frac{c^{2}}{6v}<\frac v6$, so
$\delta_{0}<\frac v6<v<A$. Put $\sigma:=\delta_{0}/A\in(0,\frac16)$. Since $n+A\ge A$ for
$n\ge0$, we have $\zeta_{j+2}\le A^{-(j-2)}\zeta_{4}$ for $j\ge2$, and hence
\[
  U_{0}\le\zeta_{4}A^{2}\sum_{j\ge3}(j+1)\sigma^{j}
   =\zeta_{4}A^{2}\,\frac{\sigma^{3}(4-3\sigma)}{(1-\sigma)^{2}} .
\]
As $\sigma<\frac16$ we have $\frac{4-3\sigma}{(1-\sigma)^{2}}\le\frac{4}{(5/6)^{2}}
=\frac{144}{25}$, and as $\delta_{0}<\frac v6$ we have $3(v-\delta_{0})>\frac{5v}{2}$.
Moreover $\delta_{0}<\frac{c^{2}}{6v}<\frac{v}{6}$, so $\delta_{0}^{2}<\frac{v^{2}}{36}$.
Therefore \eqref{eq:lbmargin} follows from
\[
  \zeta_{4}\,\frac{\delta_{0}^{2}}{A}\cdot\frac{144}{25}\;\le\;\zeta_{4}\cdot\frac{5v}{2},
  \qquad\text{which is implied by}\qquad
  \frac{144}{25}\cdot\frac{v^{2}}{36A}\;\le\;\frac{5v}{2},
\]
and the last inequality reads $\frac{4v^{2}}{25A}\le\frac{5v}{2}$, i.e.\
$\frac{8v}{25}\le5(1+v)$, which holds for every $v>0$. This proves the lower bound.

For the upper bound we estimate $W$. Since $n+A\ge A$, for $k\ge2$ and $n\ge0$ we have
$(n+A)^{-2k}\le A^{-(2k-4)}(n+A)^{-4}$, hence $\zeta(2k,A)\le A^{-(2k-4)}\zeta(4,A)$;
therefore, substituting $k=j+3$ in the sum,
\[
  W\le\zeta(4,A)\sum_{k\ge3}\frac{c^{2k-2}}{kA^{2k-4}}
   =\frac{\zeta(4,A)c^{4}}{A^{2}}\sum_{j\ge0}\frac{\rho^{2j}}{j+3}
   \le\frac{\zeta(4,A)c^{4}}{A^{2}}\cdot\frac{1}{3(1-\rho^{2})}
   =\frac{c^{2}}{2}\zeta(4,A)\cdot\frac{2\rho^{2}}{3(1-\rho^{2})},
\]
where we used $\frac{1}{j+3}\le\frac13$ and $\rho<1$. Inserting this into
\eqref{eq:balance} and dropping $U\ge0$ gives
$2\delta\zeta(3,A)\le\frac{c^{2}}{2}\zeta(4,A)\big(1+\frac{2\rho^{2}}{3(1-\rho^{2})}\big)$,
the upper bound in \eqref{eq:locbnd}.

\emph{(iii) Asymptotics.} Write $\zeta_{j}:=\zeta(j,A)$. From
$\Phi=\zeta_{2}+\frac{c^{2}}{2}\zeta_{4}+\frac{c^{4}}{3}\zeta_{6}+O(c^{6})$ and
$\zeta(2,A-\delta)=\zeta_{2}+2\delta\zeta_{3}+3\delta^{2}\zeta_{4}+O(\delta^{3})$,
equating and inserting $\delta=\alpha c^{2}+\beta c^{4}+O(c^{6})$ gives
$2\zeta_{3}\alpha=\frac{\zeta_{4}}{2}$ and
$2\zeta_{3}\beta+3\zeta_{4}\alpha^{2}=\frac{\zeta_{6}}{3}$, that is
$\alpha=\frac{\zeta_{4}}{4\zeta_{3}}$ and
$\beta=\frac{1}{2\zeta_{3}}\big[\frac{\zeta_{6}}{3}
-\frac{3\zeta_{4}^{3}}{16\zeta_{3}^{2}}\big]$, which is \eqref{eq:locasym}. \qed

\subsection{Proof of Theorem \ref{thm:sharp}}
Fix $v>0$, let $0<c<v$, $s=\sqrt{v^{2}-c^{2}}$, $A=1+v$, $\Delta=\frac{v-s}{3}$ and
$\Phi(c)=\log\Gs(x,y)/c^{2}=\sum_{k\ge1}\frac{c^{2k-2}}{k}\zeta(2k,A)$. Since
$\zeta(2,1+t)=\Phi(c)$ and $\zeta(2,\cdot)$ is strictly decreasing,
\[
  \theta>\tfrac23\iff t>\tfrac{2v+s}{3}=v-\Delta
  \iff\zeta(2,A-\Delta)>\Phi(c)\iff\mathcal{R}(c)>0 ,
\]
where $\mathcal{R}(c):=\zeta(2,A-\Delta(c))-\Phi(c)$, and $\mathcal{R}(0)=0$ by continuity.
We prove $\mathcal{R}'(c)>0$ on $(0,v)$. Because
$\Delta'=\frac{c}{3s}$ and $\partial_{\Delta}\zeta(2,A-\Delta)=2\zeta(3,A-\Delta)$,
\begin{equation}\label{eq:Rp}
  \mathcal{R}'(c)=\frac{2c}{3s}\zeta(3,A-\Delta)-\frac{2}{c}E(c),\qquad
  E(c)=\sum_{m\ge1}\frac{m}{m+1}c^{2m}\zeta(2m+2,A),
\end{equation}
so it suffices to prove
\begin{equation}\label{eq:target}
  E(c)<\frac{c^{2}}{3s}\zeta(3,A-\Delta).
\end{equation}

\emph{Step 1 (a global midpoint bound for $E$).} Put
$K_c(\xi)=\frac{1}{c^{2}}\log\frac{1}{1-c^{2}/\xi^{2}}$ and
$\eta_c(\xi)=\frac{1}{\xi^{2}-c^{2}}-K_c(\xi)$. From
$K_c(\xi)=\sum_{j\ge0}\frac{c^{2j}}{j+1}\xi^{-2j-2}$ and
$\frac{1}{\xi^{2}-c^{2}}=\sum_{j\ge0}c^{2j}\xi^{-2j-2}$,
\[
  \eta_c(\xi)=\sum_{m\ge1}\frac{m}{m+1}c^{2m}\xi^{-2m-2}>0,\qquad
  E(c)=\sum_{n\ge0}\eta_c(A+n).
\]
All coefficients being positive, $\eta_c$ is positive, decreasing and convex on $\xi>c$; the
midpoint rule for a convex function, summed over $n\ge0$, gives with $B=A-\frac12>c$
\begin{equation}\label{eq:mid}
  E(c)<\int_{B}^{\infty}\eta_c(\xi)\,d\xi ,
\end{equation}
with no restriction on the size of $c^{2}/A$. Term-by-term integration with
$y=c^{2}/B^{2}\in(0,1)$ gives
\begin{equation}\label{eq:F}
  \int_{B}^{\infty}\eta_c(\xi)\,d\xi=\frac{c^{2}}{6B^{3}}F(y),\qquad
  F(y)=6\sum_{m\ge1}\frac{m}{(m+1)(2m+1)}y^{m-1}.
\end{equation}

\emph{Step 2 (a one-variable comparison).} Let
$P(y)=\sum_{m\ge0}\frac{y^{m}}{(m+1)(2m+1)}$, so $F=6P'$, and let
$h(y)=\frac{1-\sqrt{1-y}}{2+\sqrt{1-y}}$, for which a direct differentiation gives
$6h'(y)=\frac{9}{\sqrt{1-y}\,(2+\sqrt{1-y})^{2}}$. The identity
$h=\frac{3(1-\sqrt{1-y})-y}{3+y}$, together with
$1-\sqrt{1-y}=\sum_{m\ge1}c_my^{m}$,
$c_m=\frac{\binom{2m}{m}}{4^{m}(2m-1)}$, $\frac{c_m}{c_{m-1}}=\frac{2m-3}{2m}$,
shows that $h=\sum_{m\ge1}a_my^{m}$ with $a_1=\frac16$ and
$a_m=c_m-\frac13a_{m-1}$ for $m\ge2$. Writing $P=1+\sum_{m\ge1}b_my^{m}$,
$b_m=\frac{1}{(m+1)(2m+1)}$, we have $a_1=b_1=\frac16$ and, for $r_m=a_m/c_m$,
$r_m=1-\frac{2m}{3(2m-3)}r_{m-1}$ with $r_2=\frac59$, $r_3=\frac{17}{27}$; a two-line
induction gives $\frac{3(m-1)}{4m}\le r_m<\frac34$ for every $m\ge2$. Hence
$a_2-b_2=\frac{1}{360}$, $a_3-b_3=\frac{11}{3024}$, and for $m\ge4$
$a_m>c_m-\frac14c_{m-1}=\frac{3(m-2)}{2(2m-3)}c_m$, while
$\lambda_m:=\frac{3(m-2)c_m}{2(2m-3)b_m}$ satisfies $\lambda_4=\frac{135}{128}>1$ and
$\frac{\lambda_{m+1}}{\lambda_m}=\frac{(m-1)(m+2)(2m-3)(2m+3)}{2(m-2)(m+1)^{2}(2m+1)}>1$
because $2m^{3}-5m^{2}+5m+22>0$. Thus $a_m>b_m$ for all $m\ge2$, so $h'>P'$ and
\begin{equation}\label{eq:sharp1}
  F(y)<\frac{9}{\sqrt{1-y}\,(2+\sqrt{1-y})^{2}},\qquad 0<y<1 .
\end{equation}

\emph{Step 3 (the finite-$v$ correction only helps).} Let $z=A-\Delta$,
$e=\frac{1}{2B}\in(0,1)$ and $q=\frac{s}{B}$, so that
$q=\sqrt{(1-e)^{2}-y}$ and, exactly, $z=\frac{B(2+q+4e)}{3}$. A direct computation gives
\[
  \frac{B^{3}(z+1)}{sz^{3}}=L_e(y):=\frac{9\,(2+q+10e)}{q\,(2+q+4e)^{3}},
\]
and $L_0(y)$ is precisely the right-hand side of \eqref{eq:sharp1}. Since
$q'(e)=-\frac{1-e}{q}$,
\[
  \frac{d}{de}\log L_e=-\frac{N(e,q)}{q^{2}(2+q+4e)(2+q+10e)},
\]
where
$N(e,q)=40e^{3}+40e^{2}q-12e^{2}+83eq^{2}-32eq-24e+2q^{3}+q^{2}-8q-4$. For fixed
$e\in(0,1)$ the function $q\mapsto N(e,q)$ is strictly convex
($\partial_{q}^{2}N=166e+12q+2>0$), and on $0<q<1-e$ its endpoint values are
$N(e,0)=4(e-1)(2e+1)(5e+1)<0$ and
$N(e,1-e)=9(e-1)(9e^{2}-2e+1)<0$. A convex function lies below the maximum of its
endpoint values, so $N(e,q)<0$ on $0<q<1-e$; hence $e\mapsto L_e(y)$ is strictly
increasing and
\begin{equation}\label{eq:sharp2}
  \frac{9}{\sqrt{1-y}\,(2+\sqrt{1-y})^{2}}<L_e(y)=\frac{B^{3}(z+1)}{sz^{3}} .
\end{equation}

\emph{Step 4 (conclusion).} Combining \eqref{eq:mid}--\eqref{eq:sharp2},
\[
  E(c)<\frac{c^{2}}{6B^{3}}F(y)<\frac{c^{2}}{6B^{3}}\cdot\frac{B^{3}(z+1)}{sz^{3}}
        =\frac{c^{2}(z+1)}{6sz^{3}} .
\]
Finally, the trapezoid rule applied to the convex decreasing function
$\tau\mapsto(z+\tau)^{-3}$ gives
$\zeta(3,z)\ge\int_{0}^{\infty}(z+\tau)^{-3}\,d\tau+\frac{1}{2z^{3}}=\frac{z+1}{2z^{3}}$, so
$E(c)<\frac{c^{2}}{3s}\zeta(3,z)$, which is \eqref{eq:target}. Therefore
$\mathcal{R}'(c)>0$ on $(0,v)$; with $\mathcal{R}(0)=0$ this gives $\mathcal{R}(c)>0$, and
by \eqref{eq:Rp} and the first display $\theta>\frac23$, as claimed. \qed

\section{Applications}\label{sec:app}

\subsection{Sharpened Gurland inequalities}
Since $\log\Gg(x,y)=\sum_{k\ge1}\frac{u^{2k}}{k}\zeta(2k,v)$ with all terms positive, and
since $\zeta(2k,v)\le v^{-2k}\big(1+\frac{v}{2k-1}\big)$ by the integral test, we obtain a
two-sided refinement of Gurland's inequality \cite{Gurland56}.

\begin{corollary}\label{cor:gurland}
For all $x\ne y>0$, with $v=\frac{x+y}{2}$ and $\rho_{\Gg}=\frac{|x-y|}{x+y}$,
\[
  \frac{(x-y)^{2}}{4}\,\zeta\Big(2,\frac{x+y}{2}\Big)\;\le\;\log\Gg(x,y)\;\le\;
  -\log\big(1-\rho_{\Gg}^{2}\big)+v\,\Lambda(\rho_{\Gg}),
\]
and $\Gg(x,y)>\exp\big(\frac{(x-y)^{2}}{4}\zeta(2,\frac{x+y}{2})\big)>1$. Moreover
$\log\Gg(x,y)>\frac{(x-y)^{2}}{2(x+y)}$.
\end{corollary}

\begin{proof}
The lower bound is the first term of \eqref{eq:main1}. For the upper bound, the
terms of $\zeta(2k,v)=\sum_{n\ge0}(n+v)^{-2k}$ with $n\ge1$ are dominated by the
integral over $[0,\infty)$: since $(\tau+v)^{-2k}$ decreases,
$\int_{n-1}^{n}(\tau+v)^{-2k}\,d\tau>(n+v)^{-2k}$ for every $n\ge1$, so
$\zeta(2k,v)\le v^{-2k}+\int_{0}^{\infty}(\tau+v)^{-2k}\,d\tau
=v^{-2k}\big(1+\frac{v}{2k-1}\big)$; then, with $\rho=\rho_{\Gg}$,
\[
  \log\Gg(x,y)\le\sum_{k\ge1}\frac{u^{2k}}{k}\Big(\frac{1}{v^{2k}}
   +\frac{v^{1-2k}}{2k-1}\Big)
  =\sum_{k\ge1}\frac{\rho^{2k}}{k}+v\sum_{k\ge1}\frac{\rho^{2k}}{k(2k-1)}
  =-\log\big(1-\rho^{2}\big)+v\,\Lambda(\rho).
\]
The last claim uses
$\zeta(2,v)=\sum_{n\ge0}(n+v)^{-2}>\int_{0}^{\infty}(\tau+v)^{-2}\,d\tau=\frac1v$,
because $\int_{n}^{n+1}(\tau+v)^{-2}\,d\tau<(n+v)^{-2}$ for every $n\ge0$.
\end{proof}

\subsection{Bounds for the beta function}
Recall $B(x,y)=\Gamma(x)\Gamma(y)/\Gamma(x+y)$ and the duplication formula
$\Gamma(2v)=\frac{2^{2v-1}}{\sqrt\pi}\Gamma(v)\Gamma(v+\frac12)$, where $v=\frac{x+y}{2}$.
Hence $B(x,y)=\Gg(x,y)\frac{\sqrt\pi\,\Gamma(v)}{2^{2v-1}\Gamma(v+\frac12)}$. Combining
Corollary \ref{cor:gurland} with Gautschi's inequality
$v^{1/2}<\frac{\Gamma(v+1)}{\Gamma(v+1/2)}<(v+1)^{1/2}$ \cite{Gautschi74} yields the
following.

\begin{corollary}\label{cor:beta}
For all $x,y>0$, with $v=\frac{x+y}{2}$ and $\rho_{\Gg}=\frac{|x-y|}{x+y}$,
\[
  \frac{\sqrt\pi\;v^{-1/2}}{2^{2v-1}}\,
   e^{\frac{(x-y)^{2}}{4}\zeta(2,v)}
  \;<\;B(x,y)\;<\;
  \frac{\sqrt\pi\,(v+1)^{1/2}}{2^{2v-1}\,v}\,
   e^{\,v\Lambda(\rho_{\Gg})-\log(1-\rho_{\Gg}^{2})}.
\]
\end{corollary}

\begin{proof}
Gautschi's inequality at the point $v$, with exponent $\frac12$, reads
$v^{1/2}<\frac{\Gamma(v+1)}{\Gamma(v+1/2)}<(v+1)^{1/2}$; since $\Gamma(v+1)=v\Gamma(v)$,
this gives $v^{-1/2}<\frac{\Gamma(v)}{\Gamma(v+1/2)}<\frac{(v+1)^{1/2}}{v}$. Insert this
into the expression for $B(x,y)$ and use Corollary \ref{cor:gurland}.
\end{proof}

\begin{example}
For $(x,y)=(1,2)$ one has $v=\frac32$, $\rho_{\Gg}=\frac13$, $\Gg(1,2)=1.273240\ldots$
and $B(1,2)=\frac12$. The bounds of Corollary \ref{cor:beta} read
$0.457050<0.5<0.622777$.
\end{example}

\subsection{Certified computation}
The two-sided estimate \eqref{eq:rembnd} turns the expansion into an algorithm with a
certified error.

\begin{corollary}\label{cor:alg}
Let $x\ne y>0$, $\rho=\rho(x,y)\in(0,1)$, $A=1+v$ and $0<\eps<\tfrac12$. If
\[
  m\;\ge\;\Big\lceil\frac{\log\big((1+A)/(\eps(1-\rho^{2}))\big)}{2\log(1/\rho)}\Big\rceil ,
\]
then $\big|\log\Gs(x,y)-S_m(x,y)\big|\le\eps$. Moreover, for fixed $x,y$ the number of
terms required does not exceed
\[
  \frac{\log(1/\eps)+\log(2+v)+\frac{1}{1-\rho}}{2\,(1-\rho)},
\]
so it is $O\big(\frac{\log(1/\eps)}{1-\rho}\big)$ as $\eps\to0$, the implied constant
being absolute; the arguments enter only through $\rho$ and through the additive
term $\log(2+v)$, which is independent of $\eps$.
\end{corollary}

\begin{proof}
For every $m\ge1$ we have $2m-1\ge1$, hence $1+\frac{A}{2m-1}\le1+A$ and $\frac1m\le1$;
therefore \eqref{eq:rembnd} gives
$0\le R_m\le\frac{(1+A)\rho^{2m}}{1-\rho^{2}}$, which is at most $\eps$ whenever
$\rho^{2m}\le\frac{\eps(1-\rho^{2})}{1+A}$. Taking logarithms gives the stated threshold.
For the estimate, use $\log\frac{1}{1-\rho^{2}}\le\frac{1}{1-\rho}$ (valid for
$0\le\rho<1$, since $\log\frac1{1-\rho^2}\le\frac{\rho^2}{1-\rho^2}\le\frac1{1-\rho}$) and
$-\log\rho\ge1-\rho$; then
\[
  m\le\frac{\log(1+A)+\log\frac1\eps+\log\frac{1}{1-\rho^{2}}}{2\log\frac1\rho}
  \le\frac{\log(1/\eps)+\log(2+v)+\frac{1}{1-\rho}}{2(1-\rho)} .
\]
The leading term as $\eps\to0$ is $\frac{\log(1/\eps)}{2(1-\rho)}$, with the absolute
constant $\frac12$.
\end{proof}

\subsection{The multivariate Jensen--Gurland inequality}
Theorem \ref{thm:multi} with $n=2$ reproduces Theorem \ref{thm:main}. For a concrete
three-point case take $0<\eps<\tfrac12$ and
\[
  x_{1}=x_{2}=v(1+\eps),\qquad x_{3}=v(1-2\eps),
\]
so that $\bar x=v$ and $\mu_{j}=v^{j}\eps^{j}\big(2+(-2)^{j}\big)$; \eqref{eq:multi} becomes
\[
  \log\frac{\Gamma\big(v(1+\eps)\big)^{2}\Gamma\big(v(1-2\eps)\big)}{\Gamma(v)^{3}}
   =3v^{2}\eps^{2}\zeta(2,v)+2v^{3}\eps^{3}\zeta(3,v)
    +\tfrac92v^{4}\eps^{4}\zeta(4,v)+O(\eps^{5}),
\]
all displayed coefficients being positive, in accordance with the multivariate
Jensen--Gurland inequality of Theorem \ref{thm:multi}.

\section{Numerical verification}\label{sec:num}

All computations use \texttt{mpmath} at $35$--$60$ decimal digits.
The script \texttt{verify.py} reports \texttt{ok} or \texttt{FAIL} on each of its $46$ checks.
Table \ref{tab:exp} compares the exact expansion \eqref{eq:main2} with $\log\Gs$ computed
from $\log\Gamma$. Table \ref{tab:rem} checks the two-sided estimate \eqref{eq:rembnd},
including the extreme case $(0.1,199.9)$, where the two sides differ by a factor of about
fifty. Table \ref{tab:brk} checks the bracket \eqref{eq:locbnd}; at $(0.1,199.9)$ the
upper bound equals $514.83$, against $\delta=26.6699$. Table \ref{tab:inf} traces the
approach of the ratio \eqref{eq:conj} to $\frac23$, which is reached only in the joint
limit $v\to\infty$, $\eps\to0$, $v\eps\to0$.

\begin{table}[htbp]\centering\small
\caption{Exact expansion: $|\log\Gs-S_m|$ for the truncated series \eqref{eq:main2}, and
the number $m$ of terms used.}
\label{tab:exp}
\begin{tabular}{rrrr}
\toprule
$x$ & $y$ & $|\log\Gs-S_{m}|$ & $m$\\
\midrule
$1$ & $2$ & $3.5\cdot10^{-51}$ & $34$\\
$1$ & $10$ & $5.1\cdot10^{-48}$ & $142$\\
$0.5$ & $3$ & $1.0\cdot10^{-49}$ & $68$\\
$0.2$ & $0.9$ & $2.1\cdot10^{-51}$ & $37$\\
$3$ & $3.0001$ & $1.0\cdot10^{-61}$ & $6$\\
\bottomrule
\end{tabular}
\end{table}

\begin{table}[htbp]\centering\small
\caption{The two-sided estimate \eqref{eq:rembnd}: $\mathrm{lo}=\frac{u^{2m}}{m}\zeta(2m,1+v)$
and $\mathrm{hi}$ is the right-hand side of \eqref{eq:rembnd}. The rate $\rho$ is optimal;
the ratio $\mathrm{hi}/\mathrm{lo}$ is not close to $1$ in the last row.}
\label{tab:rem}
\begin{tabular}{lrrrrr}
\toprule
$(x,y)$ & $\rho$ & $m$ & $R_m$ & lo & hi\\
\midrule
$(1,10)$ & $0.6923$ & $2$ & $4.1203\cdot10^{-1}$ & $3.1211\cdot10^{-1}$ & $6.9851\cdot10^{-1}$\\
$(1,10)$ & $0.6923$ & $12$ & $2.2659\cdot10^{-5}$ & $1.2662\cdot10^{-5}$ & $3.0165\cdot10^{-5}$\\
$(0.5,3)$ & $0.4545$ & $2$ & $3.6843\cdot10^{-2}$ & $3.2677\cdot10^{-2}$ & $5.1563\cdot10^{-2}$\\
$(0.5,3)$ & $0.4545$ & $12$ & $6.2365\cdot10^{-10}$ & $5.0458\cdot10^{-10}$ & $7.1160\cdot10^{-10}$\\
$(0.1,199.9)$ & $0.9891$ & $12$ & $2.9021$ & $0.3147$ & $15.9456$\\
\bottomrule
\end{tabular}
\end{table}

\begin{table}[htbp]\centering\small
\caption{The bracket \eqref{eq:locbnd} of Theorem \ref{thm:loc}: $\delta=v-t$ against its
lower bound $\mathrm{lo}$ and the two-term value
$\mathrm{pre}=\frac{c^{2}\zeta(4,A)}{4\zeta(3,A)}$. The upper bound is
$\mathrm{pre}\cdot\big(1+\frac{2\rho^{2}}{3(1-\rho^{2})}\big)$; at $(0.1,199.9)$ it equals
$514.83$.}
\label{tab:brk}
\begin{tabular}{lrrrr}
\toprule
$(x,y)$ & $\rho$ & $\delta=v-t$ & lo & pre\\
\midrule
$(1,2)$ & $0.2000$ & $0.0199947$ & $0.0115425$ & $0.0197486$\\
$(0.5,3)$ & $0.4545$ & $0.1186430$ & $0.0635181$ & $0.1108211$\\
$(2,7)$ & $0.4545$ & $0.2194361$ & $0.1090893$ & $0.2063217$\\
$(1,10)$ & $0.6923$ & $0.6573045$ & $0.2924374$ & $0.5586888$\\
$(20,100)$ & $0.6557$ & $4.9818888$ & $2.2129255$ & $4.4074125$\\
$(0.1,199.9)$ & $0.9891$ & $26.6699099$ & $8.2958322$ & $16.5501874$\\
$(100,101)$ & $0.0049$ & $0.0004125$ & $0.0002068$ & $0.0004125$\\
\bottomrule
\end{tabular}
\end{table}

\begin{table}[htbp]\centering\small
\caption{The ratio $\theta=\big(t-\sqrt{xy}\big)/\big(\frac{x+y}{2}-\sqrt{xy}\big)$ along
$x=v(1+\eps)$, $y=v(1-\eps)$. For a fixed $\eps$ it decreases to
$\frac23+g(\eps)$ with $g(\eps)>0$. Numerically $g(\eps)\sim\eps^{2}/180$ as $\eps\to0$
(for instance $g(10^{-1})=5.60\cdot10^{-5}$, while $10^{-2}/180=5.56\cdot10^{-5}$), so the
value $\frac23$ is approached only in the joint limit $v\to\infty$, $\eps\to0$,
$v\eps\to0$. The column $\eps=v^{-2}$ is such a path, and there
$\theta=\frac23+\frac{1}{6v}+O(v^{-2})$, as Remark \ref{rem:limit} shows.}
\label{tab:inf}
\begin{tabular}{rrrrr}
\toprule
$v$ & $\eps=10^{-1}$ & $\eps=10^{-3}$ & $\eps=v^{-2}$ & limit\\
\midrule
$10^{2}$ & $0.668396737$ & $0.668333257$ & $0.668333250$ & $0.666666667$\\
$10^{3}$ & $0.666890050$ & $0.666833339$ & $0.666833333$ & $0.666666667$\\
$10^{4}$ & $0.666739371$ & $0.666683339$ & $0.666683333$ & $0.666666667$\\
$10^{5}$ & $0.666724303$ & $0.666668339$ & $0.666668333$ & $0.666666667$\\
$10^{6}$ & $0.666722796$ & $0.666666839$ & $0.666666833$ & $0.666666667$\\
$10^{7}$ & $0.666722645$ & $0.666666689$ & $0.666666683$ & $0.666666667$\\
\bottomrule
\end{tabular}
\end{table}

Several further checks are worth recording. First, at $(x,y)=(100,1)$ one has
$Q(100,1)=4.5>1$ and $\rho(100,1)=0.96117<1$, and the expansion \eqref{eq:main2} converges,
with tails $R_{20}=1.478\cdot10^{-1}$, $R_{100}=4.359\cdot10^{-5}$ and
$R_{1000}=5.117\cdot10^{-37}$. Second, the rate in Theorem \ref{thm:rem} is confirmed: for
$(x,y)=(2,7)$ one has $2\log\rho=-1.5769147$, while the computed values of
$\frac1m\log R_m$ are $-1.8184$ ($m=5$), $-1.7829$ ($m=10$), $-1.7157$ ($m=20$),
$-1.6635$ ($m=40$), $-1.6288$ ($m=80$) and $-1.6072$ ($m=160$), approaching the predicted
limit monotonically. Third, the three-term expansion \eqref{eq:locasym} is confirmed: for
$(x,y)=(10,11)$ one has $v-t=0.00378192138106$, while the two-term formula gives
$0.00378011106367$ and the three-term formula $0.00378191965719$, the residual being
$-1.7\cdot10^{-9}$. Fourth, the large-argument asymptotics of Theorem \ref{thm:asym}: for
$(x,y)=(10^{6},2\cdot10^{6})$, where $v=1.5\cdot10^{6}$ and
$\rho=\frac{|x-y|}{2+x+y}=0.333333\ldots$, one has $\log\Gs=169898.977904\ldots$ against the
prediction $(1+v)\Lambda(\rho)=\frac{u^{2}}{1+v}\ell(\rho)=169898.919012\ldots$, a relative
error of $3.5\cdot10^{-7}$, which decreases as the arguments grow ($2.7\cdot10^{-5}$ at
$(10^{4},3\cdot10^{4})$). Fifth, Corollaries \ref{cor:gurland} and \ref{cor:beta} were
verified on the sampled configurations; for instance $B(1,2)=\frac12$ lies between
$0.457050$ and $0.622777$, and the three-point expansion displayed in
Section \ref{sec:app} agrees with $\log\Gamma$ to $19$ digits.

\section{Concluding remarks}

Theorem \ref{thm:main} is the classical Hurwitz-zeta expansion \eqref{eq:classical} of
$\log\Gamma$, specialized to the second-order symmetric difference that defines
$\log\Gs$. The poles of $\Gamma$ at the non-positive integers determine both the
coefficients and the region of convergence, which on the positive quadrant is the whole
domain, so $\log\Gs$ coincides with its series for every $x,y>0$. The condition $Q<1$ is
not required for that identification, and this settles Open Problems 1 and 2 of
\cite{Wisniewska25}. The localization of the parameter $t$ in \eqref{eq:Wt} is the content
of Theorems \ref{thm:loc} and \ref{thm:sharp}: uniqueness, the expansion of $v-t$ through
order $c^{4}$, the bracket \eqref{eq:locbnd}, and the bound $\theta>\frac23$, whose
constant is optimal and is not attained. The midpoint inequality
\eqref{eq:midpoint} is the case $\theta>\frac12$ of that bound. Two directions seem worth
pursuing:
\begin{enumerate}
\item[(a)] the $q$-analogue, replacing $\Gamma$ by the $q$-gamma function $\Gamma_{q}$,
  for which the role of $\zeta(2k,\cdot)$ is played by the Jackson $q$-zeta values, and the
  sign-regularity of the corresponding kernel;
\item[(b)] quantitative versions of Theorem \ref{thm:multi}, such as remainder estimates
  uniform in $n$ and a central limit theorem for the corresponding Jensen gaps.
\end{enumerate}

\medskip
\noindent\textbf{Statements and declarations.} The authors have no competing interests to
declare. All numerical results reported in Section \ref{sec:num} were produced by the
accompanying script \texttt{verify.py} and can be reproduced with \texttt{mpmath}
by running \texttt{python verify.py}.

\begin{thebibliography}{99}
\bibitem{Alzer97} H.~Alzer, \emph{On some inequalities for the gamma and psi functions},
Math. Comp. \textbf{66} (1997), 373--389.
\bibitem{Alzer01} H.~Alzer, \emph{Sharp inequalities for the beta function}, Indag. Math.
(N.S.) \textbf{12} (2001), 15--21.
\bibitem{Berndt85} B.~C.~Berndt, \emph{The gamma function and the Hurwitz zeta-function},
Amer. Math. Monthly \textbf{92} (1985), no.~2, 126--130,
DOI 10.1080/00029890.1985.11971552.
\bibitem{ChenChoi17} C.-P.~Chen and J.~Choi, \emph{Completely monotonic functions related
to Gurland's ratio for the gamma function}, Math. Inequal. Appl. \textbf{20} (2017),
651--659.
\bibitem{Gautschi74} W.~Gautschi, \emph{Some mean value inequalities for the gamma
function}, SIAM J. Math. Anal. \textbf{5} (1974), 282--292.
\bibitem{Gurland56} J.~Gurland, \emph{An inequality satisfied by the gamma function},
Skand. Aktuarietidskr. (1956), 171--172.
\bibitem{Kairies83} H.~Kairies, \emph{An inequality for Krull solutions of a certain
difference equation}, in: General Inequalities 3 (Oberwolfach 1981), Internat.
Schriftenreihe Numer. Math. \textbf{64}, Birkh\"auser, Basel, 1983, 277--280.
\bibitem{LaforgiaNatalini13} A.~Laforgia and P.~Natalini, \emph{On some inequalities for
the gamma function}, Adv. Dyn. Syst. Appl. \textbf{8} (2013), 261--267.
\bibitem{Lin13} L.~Lin, \emph{Further refinements of Gurland's formula for $\pi$},
J. Inequal. Appl. \textbf{2013} (2013), 2013:48.
\bibitem{Merkle05} M.~Merkle, \emph{Gurland's ratio for the gamma function},
Comput. Math. Appl. \textbf{49} (2005), 389--406.
\bibitem{MoustafaMahmoud26} H.~Moustafa and M.~Mahmoud, \emph{New bounds for Gurland-type
ratios in terms of the generalized Nielsen's beta function}, Gulf J. Math. \textbf{23}
(2026), no.~2, 1--18, DOI 10.56947/x7dwnj62.
\bibitem{DLMF} NIST Digital Library of Mathematical Functions, Release 1.2.4,
\url{https://dlmf.nist.gov/}, \S5.15 (polygamma
functions), \S25.11 (Hurwitz zeta function).
\bibitem{Qi10} F.~Qi, \emph{Bounds for the ratio of two gamma functions},
J. Inequal. Appl. \textbf{2010} (2010), Article ID 493058.
\bibitem{Srivastava88} H.~M.~Srivastava, \emph{Sums of certain series of the Riemann zeta
function}, J. Math. Anal. Appl. \textbf{134} (1988), no.~1, 129--140,
DOI 10.1016/0022-247X(88)90013-3.
\bibitem{SrivastavaChoi01} H.~M.~Srivastava and J.~Choi, \emph{Series Associated with the
Zeta and Related Functions}, Kluwer Academic Publishers, Dordrecht, 2001.
\bibitem{TianYang21} J.-F.~Tian and Z.-H.~Yang, \emph{Asymptotic expansions of Gurland's
ratio and sharp bounds for their remainders}, J. Math. Anal. Appl. \textbf{493} (2021),
124545.
\bibitem{Wisniewska25} H.~Wi\'sniewska, \emph{Some inequalities for Gurland's ratio of the
gamma functions}, arXiv:2512.07028 [math.CA], 2025.
\bibitem{YangTian24} Z.-H.~Yang and J.-F.~Tian, \emph{Complete monotonicity of the remainder
of an asymptotic expansion of the generalized Gurland's ratio}, Math. Inequal. Appl.
\textbf{27} (2024), DOI 10.7153/mia-2024-27-18.
\bibitem{YangZheng16} Z.-H.~Yang and S.-Z.~Zheng, \emph{Monotonicity of a mean related to
polygamma functions with an application}, J. Inequal. Appl. \textbf{2016} (2016),
Art.~216, DOI 10.1186/s13660-016-1155-4.
\bibitem{YangZheng19} Z.-H.~Yang and S.-Z.~Zheng, \emph{Complete monotonicity and
inequalities involving Gurland's ratios of gamma functions}, Math. Inequal. Appl.
\textbf{22} (2019), no.~1, 97--109, DOI 10.7153/mia-2019-22-07.
\end{thebibliography}

\end{document}
